Tenim una espira quadrada de 5~cm de costat. Un camp magnètic en direcció perpendicular al pla de l'espira varia en funció del temps segons l'equació $B_z(t) = B_{0_z}\cos(\omega t)$, en què $B_{0_z} = 5,0\times 10^{-6}\ \mathrm{T}$ i $\omega = 6,0\times 10^{8}\ \mathrm{rad\,s^{-1}}$.

\figura[0.4]{fig1.png}

\begin{apartat}{1}
Escriviu l'expressió del flux magnètic a través de l'espira en funció del temps i calculeu-ne el valor màxim. Indiqueu explícitament totes les unitats que intervenen en l'equació.
\end{apartat}

\begin{apartat}{1}
Escriviu l'expressió de la força electromotriu induïda a l'espira.
\end{apartat}
